\documentclass{article}
\usepackage{amsmath,amssymb}
\usepackage{xurl}
\usepackage[hidelinks]{hyperref}
\setlength{\emergencystretch}{2em}
% Shared commands for notes in docs/paper-gaps/.

\DeclareUnicodeCharacter{2080}{\ifmmode{}_{0}\else\textsubscript{0}\fi}
\DeclareUnicodeCharacter{2081}{\ifmmode{}_{1}\else\textsubscript{1}\fi}
\DeclareUnicodeCharacter{2082}{\ifmmode{}_{2}\else\textsubscript{2}\fi}
\DeclareUnicodeCharacter{2099}{\ifmmode{}_{n}\else\textsubscript{n}\fi}

\newcommand{\Lean}{\textsc{Lean}}
\ExplSyntaxOn
\NewDocumentCommand{\leanid}{m}{
  \ifmmode
    \textsf{\detokenize{#1}}
  \else
    \group_begin:
    \urlstyle{sf}
    \tl_set:Nn \l_tmpa_tl {#1}
    \tl_replace_all:Nnn \l_tmpa_tl {\_} {_}
    \exp_args:NV \path \l_tmpa_tl
    \group_end:
  \fi
}
\ExplSyntaxOff
\providecommand{\tr}{\operatorname{tr}}
\newcommand{\tnleanIssueBase}{https://github.com/LionSR/TNLean/issues/}
\newcommand{\tnleanPrBase}{https://github.com/LionSR/TNLean/pull/}
\newcommand{\tnleanDocBase}{https://sirui-lu.com/TNLean/docs/}
\newcommand{\ghissue}[1]{\href{\tnleanIssueBase#1}{GitHub issue \##1}}
\newcommand{\ghpr}[1]{\href{\tnleanPrBase#1}{PR \##1}}
\newcommand{\doclink}[1]{\href{\tnleanDocBase#1.html}{\path{#1}}}

% Dirac notation and the identity operator, as in blueprint/src/macros/common.tex.
% \providecommand keeps note-local definitions authoritative.
\usepackage{dsfont}
\providecommand{\ket}[1]{|#1\rangle}
\providecommand{\bra}[1]{\langle #1|}
\providecommand{\braket}[2]{\langle #1 | #2 \rangle}
\providecommand{\ketbra}[2]{|#1\rangle\!\langle #2|}
\providecommand{\Id}{\mathds{1}}
\providecommand{\one}{\mathds{1}}
\providecommand{\C}{\mathbb{C}}
\providecommand{\MN}[1]{M_{#1}(\C)}
\providecommand{\MD}{M_D(\C)}

% Verdict marker: \gapnote{<kind>}{<status>}. Kind is the policy
% classification (clarification, local-correction, scope-restriction,
% unfaithful, false-source, open-gap); status is open, wip (elimination
% underway), resolved, or historical. Machine-read by the site index;
% prints a verdict line.
\newcommand{\gapnote}[2]{\par\noindent\textbf{Verdict:} #1 (#2).\par}

\title{The scope of the buffered low-energy filter construction}
\author{}
\date{November 2, 2026}
\begin{document}
\maketitle
\gapnote{scope-restriction}{open}

\paragraph{Source and hypotheses.}
Arad, Landau, Vazirani and Vidick,
\cite[Theorem~t:aisgapless, lines 1027--1041]{Arad2016RigorousRG},
construct a common improving family between two nearby low-energy targets.
Their Assumption (LD), lines 281--284, fixes a positive bandwidth $\eta$.
The physical construction here retains the quantifiers
$\forall\eta_{\rm LD}>0\ \exists C_{\eta_{\rm LD}}$ before the chain,
its local dimension, its interactions and the narrower width.
Thus the hidden constants may depend on this declared fixed bandwidth.
A stronger estimate uniform in an arbitrary variable bandwidth is not
asserted. The explicit logarithmic local-dimension dependence is retained.

\paragraph{Geometry and buffers.}
The common five-region intermediate construction uses
$\ell\ge1$, $a+2\ell\le b$ and $b+2\ell\le L-1$, with native interval
$J_0=[a+\ell,b+\ell+1)$.
For a prescribed interval $J=[p,r)$ longer than $4\ell$, inward effective
cuts give $J_0\subseteq J$ with $|J\setminus J_0|\le2\ell$.
The exact regional polynomial decomposition absorbs this collar cost into
the same squared count. If $r-p\le4\ell$, actual configuration matrix
units give the required improving family directly, including the zero
target. Thus the new prescribed-interval family covers empty intervals,
short intervals and intervals meeting either chain endpoint, without an
assumed partition or fit condition. Its original targets use the same raw
Hamiltonian's canonical minimum at every scale.
The statement retains $q\ge2$; the scalar local space is a separate
elementary case. This restriction is displayed in the theorem and is not
claimed to be a counterexample to the source.
The buffered cutoffs satisfy $6w\le\eta_1\le\eta_{\rm LD}$.
The source condition $\eta_1\ge2\mu/\log n$ gives exactly this condition
under $w=\mu/(3\log n)$ for $n\ge2$. The lower original cutoff is then
$\eta_1-3w=\eta_1-\mu/\log n$, as in the source.

\paragraph{One degree and one count envelope.}
The clipping threshold is $t=\ell$, where
$\ell=\lceil C_{\eta_{\rm LD}}(1+\log m+\log(1/g))^2/g\rceil$,
$g=\min\{w,1\}$ and $m\ge q\ge2$.
The actual polynomial degree is $\ell^2$ and the count envelope is
$D=(m\ell^2)^{22\ell}$.
The same choice proves both $N\le D^2$ and $D^{12}\Delta\le10^{-5}$.
Only this numerical count base is enlarged; the physical dimension and
Hamiltonian are unchanged. Increasing an upper bound on $D$ after choosing
a filter would not justify the latter inequality and is not used.

\paragraph{Decreasing the comparison loss.}
For each integer $i\ge0$, replace only the numerical count base by
$m_i=m2^i$ and use the same fixed-bandwidth constant $C$.
The resulting scale satisfies $\ell_i\ge\ell_0+64i$, and hence the
physical truncation residual decreases by at least $2^{-4i}$.
The original-to-clipped comparisons on every prescribed interval have
loss $\nu_i=1/(200\,2^{4i})$, with
$D_i=(m_i\ell_i^2)^{22\ell_i}$ and the same-degree inequality
$D_i^{12}\Delta_i\le10^{-5}$. Long intervals retain the actual clipped
polynomial and regional factors; short intervals use actual matrix units.
The physical local dimension and interactions remain unchanged.
Thus the fixed comparison floor of the index-zero construction is not
being iterated as if it vanished. The actual images of one fixed constant-viable original upper-band
space now give arbitrarily accurate original lower-band spaces.
Their tails yield the entropy bound stated below. The sharper source
approximation rate remains separate.

\paragraph{Independent scales for sharper approximation.}
The buffered spectral comparison now permits a real threshold $t$ independent
of the integer window $\ell$, retaining the actual norm bound $3t+4\ell$.
For every fixed bandwidth one can choose
$t=C_\eta(16+\log(1/g)+\log(1/\delta))$ and an integer window bounded by
$8/g+(7C_\eta)^{1/4}g^{-1/4}(16+\log(1/g)+\log(1/\delta))^{3/4}$.
The same choice gives the squared-buffer residual budget and the required
Chebyshev suppression. These are parameter and comparison results; the
actual sharp-rank improving family, normalized fixed-cut approximation and
global sequential matrix product approximation are separate subsequent steps.

\paragraph{Remaining conclusions.}
The long-interval construction assumes no global spectral gap and supplies
the actual polynomial and tensor factors of one global clipped matrix.
The short-interval branch uses actual matrix units and does not assert a
polynomial filter with an independently supplied suppression bound.
The subsequent least-dimension argument now gives an actual
$3/200$-viable regional space for $G(\eta_{\rm LD}-v)$, with dimension
at most $\dim G(\eta_{\rm LD})D_0^3$. It uses the same canonical
targets and actual count schedule, derives a finite first count comparison,
and applies the corrected sample of size
$\lfloor p/(2D^2)\rfloor$ inside the next optimal space. Its image has
dimension strictly below the preceding minimum $p$; the companion
sampling note records the correction\footnote{\url{https://sirui-lu.com/QICLean/paper-gaps/alvv16_low_density_sampling.pdf}}.
Choosing the numerical count base $m=q\lceil e^{100}\rceil$ before
the physical construction gives
\begin{align*}
 \dim W\le\dim G(\eta_{\rm LD})\exp\!\left[
 \frac{K_{\eta_{\rm LD}}}{v}
 \bigl(1+\log q+\log(1+1/v)\bigr)^3\right],
\end{align*}
where $K_{\eta_{\rm LD}}$ precedes the chain, local dimension,
interactions, interval and deficit. For $0<v\le\eta_{\rm LD}$,
the cap satisfies $v/2\le\min(v,\eta_{\rm LD}/2)$; larger deficits
give the zero target. Thus the logarithmic count estimate introduces
no additional upper bound on the requested positive deficit.
The accurate interval families now give, for every unit vector
$\psi\in G(\eta_{\rm LD}-v)$, the actual regional entropy bound
\begin{align*}
 S(\psi_J)\le\log\dim G(\eta_{\rm LD})+\frac{K_{\eta_{\rm LD}}}{v}
 \bigl(1+\log q+\log(1+1/v)\bigr)^3.
\end{align*}
The coefficient on the full-band logarithmic dimension is exactly one.
The proof fixes one viable outer-band space before all accuracy choices,
then controls tails of one actual reduced state by its image spaces.
No spectral gap or invariance of an original band under a clipped polynomial
is assumed. Every positive deficit and every interval is included.
The sharper accuracy-dependent MPS approximation rate, a low-energy basis
of matrix product states and a running-time bound remain separate. Closed spectral
cutoffs and the absence of a positivity claim for the Chebyshev polynomial
follow the companion spectral-projector and positivity notes.

\bibliographystyle{plain}
\bibliography{../../blueprint/src/references}
\end{document}
